% English translation of capitulos/04e-exame.tex, made by Claude Code (Claude Fable 5.2) on 24 November 2026.
% Pending review by the author. When the Portuguese file changes, update this one (scripts/verificar_traducao.py).
\subsection{Examination of the table in the scenarios}\label{sec:exame}
Table~\ref{tab:cenarios} summarizes the eight synthetic scenarios, two per task. Each one fixes language and domain,
model, verification resources and call budget. The evidence expectations are hypotheses to be confronted. The
scenarios and decision-support table v1 were frozen with SHA-256 before the examination. For each scenario, the examination applies the
rows of the table, checks the studies cited in the records and registers limitations and ambiguities. Claude Code drafted the examination and a proposed judgement, and the author made the final judgement. The examination does not measure effectiveness.

\begin{table}[!t]
\caption{Scenarios used in the examination of version 1 of the decision-support table}\label{tab:cenarios}
\centering\scriptsize
\setlength{\tabcolsep}{3pt}
\begin{tabular}{@{}p{0.55cm}p{3.2cm}p{2.4cm}p{1.3cm}@{}}
\hline
ID & Situation & Verification & Budget\\
\hline
G1 & Python functions from a docstring; large model & Public unit tests & 5 calls\\
G2 & Solidity 0.8 contracts; local model up to 8B & \texttt{solc} and static analyser & 3 calls\\
D1 & Defect in a Python function & Failing test and suite & 10 calls\\
D2 & Defect in a Java repository from a report & Compilation; suite does not reproduce & 20 calls\\
T1 & C99 module to safe Rust & \texttt{rustc} and ported tests & 5 rounds\\
T2 & Batch COBOL to Java 17 & Java compiler only & 3 calls\\
R1 & Long Python methods, lower complexity & Tests and complexity meter & 3 calls\\
R2 & API migration in about 200 Java files & Compilation, tests, benchmarks, changelog & 2 per file\\
\hline
\end{tabular}
\end{table}

The author judged the eight scenarios on 17 November 2026 (Table~\ref{tab:julgamento}). On the three criteria, the result was:
\begin{itemize}
  \item clarity: 1 meets and 7 partial;
  \item grounding in the sources: 2 meet, 5 partial and 1 does not meet;
  \item concrete guidance: 1 meets, 6 partial and 1 does not meet.
\end{itemize}
The author agreed with 22 of the 24 judgements proposed by Claude Code and downgraded to partial the grounding of
T1 and R1. In T1, the result cited (S732-C1) comes from a package with several repair steps and limits per step, which
is not equivalent to five total rounds per function. In R1, the highlighted results do not support at the same time the
reduction of complexity and that of smells.

\begin{table}[!t]
\caption{Author's judgement per scenario and criterion}\label{tab:julgamento}
\centering\scriptsize
\begin{tabular}{@{}lccc@{}}
\hline
Scenario & Clarity & Grounding & Guidance\\
\hline
G1 & partial & meets & partial\\
G2 & partial & partial & partial\\
D1 & partial & meets & meets\\
D2 & partial & partial & does not meet\\
T1 & meets & partial & partial\\
T2 & partial & does not meet & partial\\
R1 & partial & partial & partial\\
R2 & partial & partial & partial\\
\hline
\end{tabular}
\end{table}

D1 is the only scenario in which grounding and guidance meet the criteria: the table leads to an executable script, with the
function, the test and the error message in the prompt and verification by the suite. The guidance fails in D2, in which the
suite does not reproduce the defect and the localization is unknown, and the grounding fails in T2, in which there is only a compiler
for a COBOL migration that requires equivalence of behaviour. In both cases, the table does not declare that the
evidence is insufficient. The ambiguities that the author recorded recur across the scenarios:
\begin{itemize}
  \item the table does not say how to split a fixed budget of calls between repair and new candidates (G1, G2 and T1);
  \item the feedback condition puts together compiler, tests and static analyser and does not distinguish the result
  measured, such as compilation, passing the tests or fixing the defect (G2, T1 and T2);
  \item passing the suite does not guarantee that the change fixes the intended defect or that it preserves the behaviour
  (D1, D2 and T2);
  \item the dependence on known localization or on a test that reproduces the defect only appears in the caveats (D2);
  \item the refactoring row puts together specification, domain knowledge and examples, which have different foundations
  (R1 and R2).
\end{itemize}
The author also corrected the derived guidance in some scenarios. In G2, the appropriate wording is not to recommend
self-critique without execution as a default, since one deterioration with a 1.5B model does not establish the behaviour of all
models up to 8B. In D2, compilation and existing tests still give feedback on the change, although they do not verify
the reported defect. In T1, the study cited does not have a budget equal to that of the scenario. In R1, from one to three examples is not
a range established by the evidence.

The author decided to apply the revisions in a version 2 of the table:
\begin{itemize}
  \item distinguish the types of feedback and make explicit the result measured (compilation, passing the tests,
  correctness of the repair, complexity or performance);
  \item make explicit the dependences on localization and on reproduction of the defect;
  \item declare the conditions without sufficient evidence and the cases in which the corpus does not guide the distribution of the
  calls;
  \item separate the ingredients grouped in the refactoring row;
  \item differentiate absence of evidence, contrary result and indirect evidence.
\end{itemize}
In the first proposal, the eight hypotheses of the scenarios were given as confirmed, H-G1 only in part. The author asked
that each hypothesis be judged separately, admitting partial confirmation or an inconclusive result. The examination assesses the
guidance offered by the table, and not the effectiveness of the techniques. The complete examination is in
\path{evidencias/sinteses/exame-cenarios-v1.md}, and the author's judgements, in
\path{evidencias/sinteses/exame/julgamentos-autor.json}.
Version 2 of the table, revised after the examination, is in Section~\ref{sec:qp3}. In the separate judgement, the author confirmed six hypotheses and confirmed two in part
(\path{evidencias/sinteses/exame/hipoteses-proposta.md}). H-G1 is confirmed in part because there is evidence to
consider tests, feedback and planning, but not to choose between them with five calls. H-T1 is
confirmed in part because the result cited comes from a package with its own steps and limits and does not cover the requirement
of safe Rust. No hypothesis was inconclusive. The six confirmed include the four that predicted insufficient
evidence (H-G2, H-D2, H-T2 and H-R2): in these scenarios, not recommending a technique as a default is the appropriate
answer.
